\section{Heterogeneous complete-assignment proofs}\label{app:heterogeneous}
\begin{proof}[Proof of Theorem~\ref{thm:sharedprojection}]
Complete assignments are precisely partial matchings in the declared graph
$\Gamma$, with every unmatched service retaining a separate bus. For an
allowed matching $M$, a bus serving A$_i$ has inventory $B-a_i$ before
early charging. Headroom gives $0\le x_i\le a_i$; pairing B$_j$
adds $x_i\ge[a_i+b_j-B]_+\le a_i$. An unmatched B bus is full
throughout the early window. Summing these independent intervals gives
$[R(M),\sum_i a_i]$. All A buses are present for the entire early
window. At $K>0$ serial charging takes $E/K$ hours, so exactly
$E\le K$ is required. When $K=0$, only zero early energy is possible.

Every bus is present at hour 3. Its terminal demand is $a_i+b_j-x_i$
if paired, $a_i-x_i$ if A-only, or $b_j$ if B-only. These quantities
are nonnegative, individually no greater than $B$, and total $S-E$.
At $H>0$ serial replenishment takes $(S-E)/H$ hours, giving precisely
$S-E\le H$. At $H=0$ no positive terminal demand can be met.
Monotone charging preserves inventory bounds and restores each bus full.
Thus each matching has exactly interval $[\max\{L,R(M)\},U]$ when
nonempty. Finite minima and a common upper endpoint give the stated
cardinality union for $0\le m\le\nu$. Every such cardinality exists by
deleting edges from a maximum allowed matching. Since the empty matching
and every singleton option are retained, feasibility is exactly $L\le U$,
also when $\nu=0$ or some vertices are isolated.

Deleting an edge from a minimum size-$m$ matching gives
$r_{m-1}\le r_m$. For discrete convexity, use the standard unit-capacity
source--A--B--sink network containing only the allowed A--B edges, with
their costs $r_{ij}$. Its fixed-flow LP is integral at integer flow values
$m=0,\ldots,\nu$. The LP value is convex in required flow by duality,
so $2r_m\le r_{m-1}+r_{m+1}$ for $1\le m<\nu$. This reasoning
uses bipartiteness, not completeness or the sorted-energy shortcut.
All thresholds at or below $L$ are dominated
in cost by $m_0$. The first positive retained horizontal increment is
$r_{m_0+1}-L$, no greater than $r_{m_0+1}-r_{m_0}$; later increments
are the original nondecreasing matching increments. Hence frontier slopes
$-f/\text{increment}$ are nondecreasing, followed by zero. Endpoints are
physical and their interpolants are hull mixtures. Every physical branch
lies above this convex graph, so every mixture does too. The graph is
therefore exactly the lower hull boundary. Cardinalities above $m_{\max}$
have empty physical intervals and supply no endpoint. Cardinalities above
$\nu$ supply no matching at all. If $m_0=m_{\max}$, including the
empty-graph case, only the physical horizontal segment remains.
\end{proof}

\begin{proof}[Proof of Corollary~\ref{cor:sharedsupport}]
The convex frontier plus strictly convex continuous $G$ has one minimizing
load. At $L$, at a retained threshold or on the final horizontal segment,
a physical plan attains the frontier. Strictly between consecutive retained
thresholds, its operating cost is strictly below that of every physical
branch because $f>0$. No physical plan attains the hull minimizer; compactness
of the finite union of physical segments gives $D>\CH$. This retains
physical thresholds even on a larger collinear hull face. The mathematical
value equivalence becomes a price-support equivalence only with the
framework's separate qualifications.
\end{proof}

\begin{proof}[Proof of Corollary~\ref{cor:sortedmatching}]
Under the corollary's complete-graph assumption, replacing a selected service
by an unused smaller energy while keeping its
partner cannot increase the coordinatewise nondecreasing edge cost. Thus
some optimal size-$m$ matching uses the $m$ smallest energies in each batch.
Let $h(t)=[t-B]_+$. For $a\le a'$ and $b\le b'$, monotonicity of
the convex increment $h(t+a'-a)-h(t)$ gives
\[
 h(a+b)+h(a'+b')\ge h(a+b')+h(a'+b).
\]
Swapping any two partners ordered in the same increasing direction therefore
cannot raise cost. Repeating exchanges yields opposite order and the stated
sum. Ties and zero edge costs require no separate assumption.
\end{proof}

\begin{proof}[Proof of Theorem~\ref{thm:batchprojection}]
Every bus serves at most one A and one B, so service assignments are exactly
partial bipartite matchings. For a fixed matching $M$, let $x_i$ be the
early purchase on the bus serving $A_i$. After that service it has inventory
$20-a_i$. Capacity and early acceptance imply $0\le x_i\le u_i$.
If matched to $B_j$, its inventory after B is $20-a_i+x_i-b_j$, giving
$x_i\ge r_{ij}$. The assumptions imply $r_{ij}\le a_i$ and
$r_{ij}\le10$, hence each interval is nonempty. An unmatched B bus
starts full and has no earlier consumption, so it cannot accept early
charge. The independent early intervals sum to exactly
\[
 \left[\sum_{(i,j)\in M}r_{ij},U\right].
\]
There are $n$ A buses and $n$ early connectors, with aggregate grid
power $10n$, so every choice of these purchases is realizable.

For terminal charging, a matched bus needs $a_i+b_j-x_i\le20$ kWh
and fits alone on one connector. An unmatched A bus needs $a_i-x_i$;
an unmatched B bus needs $b_j$. Group unmatched A and B buses arbitrarily;
each group's sum is at most $a_i+b_j\le30$ and fits serially on one
30 kW connector during the terminal hour. There are $m$ matched groups
and $n-m$ unmatched groups. Aggregate power $30n$ allows all groups
simultaneously. Every bus reaches its own full inventory without intermediate
overcharge or energy deficit. Full replenishment fixes total purchase $S$.

All size-$m$ matchings have the same operating cost and upper early endpoint
$U$, and there are finitely many. Their interval union is $[r_m,U]$,
proving the exact projection without restricting the matching set.

Deleting one edge from a minimum size-$m$ matching proves
$r_{m-1}\le r_m$. To prove discrete convexity, represent the bipartite
matching as unit-capacity flow from a source through A and B vertices to
a sink, charging edge cost $r_{ij}$. At integer required flow the LP is
integral and has value $r_m$. Its value as a function of required flow is
convex by LP duality, a supremum of affine functions. Therefore
$2r_m\le r_{m-1}+r_{m+1}$.

Nonnegative nondecreasing increments can vanish only in an initial block.
After $m_0$, set $d_j=r_j-r_{j-1}>0$. The interpolant slopes
$-f/d_j$ are nondecreasing, and the final slope zero is no smaller.
The graph is convex and no higher than any physical horizontal branch;
the earlier zero-threshold branches have strictly larger cost because $f>0$.
Each interpolant edge is a convex combination of physical endpoints, and
the final horizontal segment is physical. Thus its graph belongs to the
hull. Conversely every physical point lies on or above this convex graph,
and therefore so does every convex combination. This proves it is exactly
the hull's lowest operating-cost boundary.
\end{proof}

\begin{proof}[Proof of Theorem~\ref{thm:batchsupport} and the value formulas]
The hull boundary is convex; $G$ is strictly convex with derivative
$G'(E)=2(E-E_0)/(5n)$. Their sum has a unique minimizing early load.
At a threshold or on the final horizontal segment, a physical plan attains
the hull boundary, so $D=\CH$. In an adjacent open threshold interval,
the hull operating cost is strictly below the cheapest physical branch
$c_{j-1}$, since $f>0$. No physical plan attains the unique hull-optimal
projected point, and every other load has larger hull objective. The
finite union of compact physical segments attains its minimum, proving
$D>\CH$. Proposition~\ref{prop:support} supplies the support equivalence.

At an intermediate threshold $r_m$, the objective's one-sided slopes are
$G'(r_m)-f/d_m$ and $G'(r_m)-f/d_{m+1}$. Their signs straddle zero
exactly under \eqref{eq:batchvertex}, including weak inequalities.
Strict convexity on each physical branch gives \eqref{eq:batchD}.
On a sloping hull edge its derivative is $G'(E)-f/d_j$, whose zero is
$E_0+5nf/(2d_j)$. Clamping to the edge gives \eqref{eq:batchh}; the
remaining horizontal segment minimizes $c_n+G$ directly. This covers
the $m_0=n$ case without dividing by a zero threshold increment.
\end{proof}

\begin{proof}[Details of the heterogeneous example]
The stated matrix gives $r_0=r_1=0$ and $r_2=9$. Physical minima for
$m=0,1$ occur at $E=9/2$ and cost $5841/40$ and $5561/40$,
respectively. For $m=2$, the minimum is the endpoint $E=9$, with value
$2681/20$. The sloping hull edge objective is
\[
 \frac{2821}{20}-\frac{151}{90}E+\frac{E^2}{10},
\]
whose minimum is $434201/3240$ at $151/18$. Subtracting gives
$\Delta=121/3240$. The endpoint weights stated in the text realize the
same load and operating interpolation. At $p=(49/10,4)$, the price
difference is positive, so each branch minimizes its bill at its lowest
early endpoint. Their bills plus operating costs are 224, 217 and
$2181/10$. Hence the full-set regret is $11/10$, not merely a gain
against an arbitrarily selected alternative. All calculations are exact
algebra on the displayed ideal model.
\end{proof}
